\section{Architecture}
\label{sec:Architecture}

This chapter describes the overall system architecture, component layout and
major data flows in Archimesh.

\subsection{System purpose}

The system lets multiple Archi clients collaborate on the same model timeline
through a centralized server. It combines:

\begin{itemize}
\item a local Archi client plugin that captures and applies model changes
\item a Quarkus server that validates, orders, persists and broadcasts changes
\item Kafka for ordered model event fan-out
\item Neo4j for the materialized model graph, op-log, tags and merge metadata
\end{itemize}

The model timeline is intentionally linear: one \texttt{HEAD}, no branches,
immutable tags as named historical points.

\subsection{High-level flow}

\begin{enumerate}
\item A user edits a model in Archi.
\item The client plugin captures local EMF changes and converts them into
  Archimesh operations.
\item The server validates the op batch against preconditions, locks and
  idempotency state, then assigns authoritative revisions.
\item The server persists the batch, updates the materialized graph and
  broadcasts it via Kafka.
\item Other connected clients apply the accepted operations into their local
  Archi model.
\end{enumerate}

\subsection{Component layout}

\begin{longtable}{p{38mm}p{116mm}}
\toprule
Component & Responsibility \\
\midrule
\filepath{server/} & Quarkus server: WebSocket and REST endpoints, authorization,
  validation, revision assignment, idempotency, persistence, broadcast. \\
\filepath{client/org.gautelis.archimesh.plugin/} & Eclipse RCP plugin:
  WebSocket client, EMF change capture, operation mapping, remote apply,
  offline outbox, cache persistence. \\
\filepath{client/client-tests/} & Client plugin tests: CRDT merge, operation
  mapping, folder sync, outbox replay, notation inventory. \\
\filepath{admin-web/} & Svelte admin UI: model overview, catalog, versions,
  access control, sessions, audit. \\
\filepath{schemas/} & JSON schemas for common types, op batches, operations,
  lock events, presence and notation contracts. \\
\filepath{neo4j/} & Cypher schema for constraints and indexes, plus a
  documented graph model. \\
\filepath{scripts/} & Build, validation, development and operational scripts. \\
\filepath{docker-compose.yml} & Local dev infrastructure: Neo4j + Kafka. \\
\bottomrule
\end{longtable}

\subsection{Server component map}

The server is organized into distinct packages:

\begin{longtable}{p{42mm}p{112mm}}
\toprule
Package & Responsibility \\
\midrule
\texttt{endpoint} & WebSocket endpoint (\texttt{ArchimeshEndpoint}), REST
  endpoints (\texttt{ModelStateEndpoint}, \texttt{AdminEndpoint},
  \texttt{AdminUiEndpoint}, \texttt{ConsistencyEndpoint}). \\
\texttt{service} & Central coordination (\texttt{ArchimeshService}), plus
  service interfaces for validation, revision, locks, Neo4j repository,
  Kafka publish/consume, session registry and idempotency. \\
\texttt{service.impl} & Concrete implementations: in-memory services for
  validation, revision, locks and sessions; Neo4j repository with split
  support classes for op-log, materialized state, reads and compaction;
  Kafka publisher and consumer implementations. \\
\texttt{auth} & Authorization: identity modes, role normalization,
  policy decision point (PDP), authorization service and transport. \\
\texttt{model} & Admin data classes: status, catalog, sessions, ACLs,
  export/import, compaction, integrity, audit. \\
\texttt{wire} & Wire protocol envelopes (\texttt{ClientEnvelope},
  \texttt{ServerEnvelope}) and message types for inbound
  (Join, SubmitOps, AcquireLock, ReleaseLock, Presence) and outbound
  (CheckoutSnapshot, CheckoutDelta, OpsAccepted, OpsBroadcast,
  LockEvent, PresenceBroadcast, Error). \\
\bottomrule
\end{longtable}

\subsection{Client plugin component map}

\begin{longtable}{p{42mm}p{112mm}}
\toprule
Package & Responsibility \\
\midrule
\texttt{ws} & WebSocket session manager, outbox handling, join/rejoin
  decisions, local submit tracking, inbound message dispatch. \\
\texttt{emf} & EMF change capture, operation mapping, remote op
  application, CRDT merge (LWW, OR-set, tombstones), model controller. \\
\texttt{startup} & Eclipse startup hooks, workbench lifecycle bridge,
  status line bridge. \\
\texttt{handlers} & Menu command handlers: connect, disconnect, open
  server model, switch active model, show last conflict, resynchronize. \\
\texttt{ui} & Connect dialog, server model catalog dialog, open server
  model dialog. \\
\texttt{notation} & Notation serializer and deserializer for Archi
  diagram fields. \\
\texttt{util} & Authentication hint helpers and lightweight JSON
  utilities. \\
\bottomrule
\end{longtable}

\subsection{Two kinds of state}

The system maintains two kinds of state at all times:

\begin{description}
\item[Local working state] Inside each Archi client. This is the EMF model
  the user sees and edits. It is not authoritative.
\item[Authoritative shared state] On the Archimesh server. The Neo4j graph
  is the source of truth. The op-log is the replayable history.
\end{description}

The client plugin synchronizes between these two states: capturing local
changes, submitting them to the server, and applying remote changes into the
local model.

\subsection{Snapshot vs op-log}

The server keeps both a materialized snapshot-like graph and an append-only
op-log in Neo4j. Why both exist:

\begin{itemize}
\item The materialized graph makes joins, exports and admin inspection fast
  (no replay needed).
\item The op-log preserves the authoritative committed history for rebuild,
  import and time-travel.
\end{itemize}

This is why export packages contain both the current snapshot and the full
op batch history. The snapshot is derived from the op-log; the op-log is
the ground truth.

\subsection{Folder model}

Model-tree folders are first-class synchronized state with strict rules:

\begin{itemize}
\item Root folders have stable synthetic ids (e.g. \texttt{folder:root-business}).
\item User folders have immutable ids --- rename is allowed, but identity is
  stable.
\item Object placement in folders is explicit and persisted.
\item Cross-root folder moves are rejected by validation.
\item Non-empty folder delete is rejected.
\end{itemize}

This avoids drift between clients and preserves Archi model-tree structure.

\subsection{Linear timeline}

Every model has exactly one linear revision timeline with:

\begin{itemize}
\item One moving \texttt{HEAD}
\item Zero or more immutable tags
\item No branches
\end{itemize}

Clients submit batches against a \texttt{baseRevision}. The server assigns an
authoritative contiguous revision range. If a batch contains \(N\) ops it
consumes \(N\) revisions. Tags are read-only pointers to historical revisions,
stored with a snapshot of the model at that point.

Only \texttt{HEAD} is writable. Tagged revisions are read-only for both reads
and writes.

\subsection{Deployment and operational considerations}

\subsubsection{Transport security}

The WebSocket and REST endpoints do not terminate TLS themselves. In production,
deploy behind a reverse proxy (Nginx, Caddy, Traefik) that terminates TLS and
forwards to the Quarkus server. Example configurations are provided in
\filepath{server/examples/}.

When using \texttt{app.identity.mode=proxy}, the reverse proxy must also
forward trusted identity headers, and the WebSocket upgrade path must carry
those headers through the proxy to the server.

\subsubsection{Reconnection backpressure}

After a server restart, all previously connected clients may attempt to
reconnect simultaneously. To avoid overwhelming the server:

\begin{itemize}
\item Clients should apply jitter to reconnection attempts (randomized delay
  1--5 seconds before the first retry).
\item The server should reject joins that arrive faster than a configurable
  rate limit.
\item Admin diagnostics expose active session counts so operators can detect
  reconnection storms.
\end{itemize}

The current client implementation does not yet include reconnect jitter. This
is a known operational gap for production deployments with many concurrent
clients.

\subsubsection{Kafka single-partition requirement}

Each model's ops topic must have exactly one partition to preserve total
order. If multiple partitions are used, the total order guarantee breaks and
clients can receive operations in different orders, producing divergent state.
This is a hard constraint, not a recommendation.

If Kafka publish fails, the server falls back to broadcasting
\texttt{OpsBroadcast} directly to connected WebSocket sessions, ensuring
clients on the same server instance still receive updates.
